\section{Fairness Notions}
\label{sec:notions}

We now describe all the different fairness notions we consider in this paper.

\subsection{Basic Envy-Based Notions}

We begin with the most well-known notion of fairness: envy-freeness (EF).

\begin{definition}[EF]
\label{defn:ef}
Let $\Ical \defeq (N, M, (v_i)_{i \in N}, w)$ be a fair division instance.
In an allocation $A$, an agent $i \in N$ \emph{envies} another agent $j \in N \setminus \{i\}$ if
\[ \frac{v_i(A_i)}{w_i} < \frac{v_i(A_j)}{w_j}. \]
Agent $i$ is \emph{envy-free} in $A$ (or $A$ is EF-fair to $i$) if
she doesn't envy any other agent in $A$.
\end{definition}

When entitlements are unequal, most papers use the term WEF instead of EF.
But we use the term EF in this paper to emphasize that unequal entitlements
is a property of the fair division setting, not the fairness notion.

It is easy to see that EF allocations may not exist, so several relaxtions have been studied.
Two of the most popular relaxations of EF are EF1 and EFX.
EF1 was proposed by \cite{budish2011combinatorial,lipton2004approximately},
and EFX was proposed by \cite{caragiannis2019unreasonable}.

\begin{definition}
\label{defn:ef1}
For a fair division instance $\Ical \defeq (N, M, (v_i)_{i \in N}, w)$,
an allocation $A$ is EF1-fair to agent $i$ if for every other agent $j$,
at least one of the following hold:
\begin{enumerate}
\item $i$ does not envy $j$.
\item $\exists g \in A_j$ such that
    \[ \frac{v_i(A_i)}{w_i} \ge \frac{v_i(A_j \setminus \{g\})}{w_j}. \]
\item $\exists c \in A_i$ such that
    \[ \frac{v_i(A_i \setminus \{c\})}{w_i} \ge \frac{v_i(A_j)}{w_j}. \]
\end{enumerate}
\end{definition}

\begin{definition}
\label{defn:efx}
For a fair division instance $\Ical \defeq (N, M, (v_i)_{i \in N}, w)$,
an allocation $A$ is EFX-fair to agent $i$ if for every other agent $j$,
$i$ doesn't envy $j$, or both of the following hold:
\begin{enumerate}
\item $\displaystyle \frac{v_i(A_i)}{w_i} \ge \frac{\max(\{v_i(A_j \setminus S) \mid S \subseteq A_j
    \textrm{ and } v_i(S \mid A_i) > 0 \textrm{ and } v_i(S \mid A_j \setminus S) \ge 0\})}{w_j}$.
\item $\displaystyle \frac{\min(\{v_i(A_i \setminus S) \mid S \subseteq A_i
    \textrm{ and } v_i(S \mid A_i \setminus S) < 0 \})}{w_i} \ge \frac{v_i(A_j)}{w_j}$.
\end{enumerate}
\end{definition}

\Cref{defn:efx} looks very different from the original definition of EFX
given by \cite{caragiannis2019unreasonable}.
However, in \cref{sec:notions:efx}, we explain why \cref{defn:efx} makes sense,
and we show that it's equivalant to the original definition when
valuations are submodular and the items are goods.

\subsection{Basic Share-Based Notions}

\begin{enumerate}
\item Maximin Share (MMS). For unequal entitlements, we consider weighted MMS (WMMS), which generalizes MMS.
\item AnyPrice Share (APS).
\item Proportionality (PROP).
\item Proportional up to one item (PROP1).
\item Proportional up to any item (PROPX).
\item Proportional up to maximin item (PROPM).
\end{enumerate}

\subsection{Derived Notions}

New fairness notions can be obtained by systematically modifying existing notions.

We start with two related concepts, epitemic fairness \cite{aziz2018knowledge,caragiannis2023new}
and minimum fair share \cite{caragiannis2023new}.

\begin{definition}[epistemic fairness]
\label{defn:epistemic}
Let $F$ be a fairness notion.
An allocation $A$ is \emph{epistemic-$F$-fair} to an agent $i$ if
there exists another allocation $B$ that is $F$-fair to agent $i$ and $B_i = A_i$.
$B$ is called agent $i$'s \emph{epistemic-$F$-certificate} for $A$.
\end{definition}

\begin{definition}[minimum fair share]
\label{defn:minfs}
For a fair division instance $\Ical \defeq (N, M, (v_i)_{i \in N}, w)$
and fairness notion $F$, let $\Acal(\Ical, F, i)$ be the set of allocations
that are $F$-fair to agent $i$.
Then $i$'s \emph{minimum $F$ share} is given by
\[ \minFS(\Ical, F, i) \defeq \min_{A \in \Acal(\Ical, F, i)} v_i(A_i). \]
An allocation $A$ is \emph{minimum-$F$-share-fair} to agent $i$ if
$v_i(A_i) \ge \minFS(\Ical, F, i)$.

Equivalently, $A$ is minimum-$F$-share-fair to agent $i$ if there exists
an allocation $B \in \Acal(\Ical, F, i)$ such that $v_i(A_i) \ge v_i(B_i)$.
Then $B$ is called agent $i$'s \emph{minimum-$F$-share-certificate} for $A$.
\end{definition}

We now describe pairwise fairness \cite{caragiannis2019unreasonable}
and groupwise fairness \cite{barman2018groupwise}.

\begin{definition}[restricting]
\label{defn:restricting}
Let $\Ical \defeq (N, M, (v_i)_{i \in N}, w)$ be a fair division instance
and $A$ be an allocation. For a subset $S \subseteq N$ of agents, where $|S| \ge 2$,
let $\restrict(\Ical, A, S)$ be the pair $(\Icalhat, \Ahat)$, where
$\Ahat \defeq (A_j)_{j \in S}$, $\Icalhat \defeq (S, \Mhat, (v_j)_{j \in S}, \what)$,
$\Mhat \defeq \bigcup_{j \in S} A_j$, and $\what_j \defeq w_j / \sum_{j \in S} w_j$.
\end{definition}

\begin{definition}[pairwise fairness]
\label{defn:pairwise}
For a fair diviion instance $\Ical \defeq (N, M, (v_i)_{i \in N}, w)$,
an allocation $A$ is called \emph{pairwise-$F$-fair} to agent $i$ if
for all $j \in N \setminus \{i\}$, $A^{(j)}$ is $F$-fair to $i$ in the instance $\Ical^{(j)}$,
where $(\Ical^{(j)}, A^{(j)}) \defeq \restrict(\Ical, A, \{i, j\})$.
\end{definition}

\begin{definition}[groupwise fairness]
\label{defn:groupwise}
For a fair diviion instance $\Ical \defeq (N, M, (v_i)_{i \in N}, w)$,
an allocation $A$ is called \emph{groupwise-$F$-fair} to agent $i$ if
for all $S \subseteq N \setminus \{i\}$, $A^{(S)}$ is $F$-fair to $i$ in the instance $\Ical^{(S)}$,
where $(\Ical^{(S)}, A^{(S)}) \defeq \restrict(\Ical, A, \{i\} \cup S)$.
\end{definition}

In this paper, we consider the following derivative notions:
\begin{enumerate}
\item Epistemic envy-freeness (EEF).
\item Epistemic EFX (EEFX).
\item Epistemic EF1 (EEF1).
\item Minimum EFX share (MXS).
\item Minimum EF1 share (M1S).
\item Pairwise proportionality (PPROP) and groupwise proportionality (GPROP).
\item Pairwise MMS (PMMS) and groupwise MMS (GMMS).
\item Pairwise APS (PAPS) and groupwise APS (GAPS).
\end{enumerate}
